\slajdid{02-s050}
\begin{frame}[label=02-s050]{Pojava lažnih nula i jedinica}
\gusttekst
\centering (race hazard)\par
\raggedright
Zbog kašnjenja koja su različita kod logičkih kola i zavisna i od fabrikacije i od temperature i od starosti komponente \ldots\ mogu se desiti neželjeni efekti na izlazu prilikom promene ulaznih signala. Najčešća situacija je da je izlaz na logičkoj jedinici, da se ulazni signali promene tako da izlaz treba da ostane na logičkoj jedinici ali se zbog propagacije signala pojavi kratkotrajna, lažna, nula na izlazu. Isto takva situacija i da je izlaz na logičkoj nuli, da se ulazni signali promene tako da izlaz treba da ostane na logičkoj nuli ali se zbog propagacije signala pojavi kratkotrajna, lažna, jedinica na izlazu. Dosta čest termin za ove kratkotrajne neželjene signale je glič (glitch). Da vidimo na primeru. Neka je realizovana kombinaciona mreža
\[
F=C\bar B+BA
\]
\centering
\begin{tikzpicture}[circuit logic US,DEoznaka,x=1cm,y=1cm,
 gate/.style={draw,logic gate inputs=nn,minimum width=12mm,minimum height=9mm,scale=.6}]
\node[and gate US,gate] (u) at (2,1.1) {};
\node[and gate US,gate] (d) at (2,-.1) {};
\node[or gate US,gate] (f) at (3.8,.5) {};
\node[not gate US,draw,minimum width=8mm,minimum height=8mm,scale=.6] (i) at (.6,.55) {};
\draw[DEvod] (u.input 1)--(-.5,0 |- u.input 1) node[left] {$C$};
\draw[DEvod] (i.input)--(-.5,.55) node[left] {$B$};
\draw[DEvod] (i.output)--(1.25,.55)|-(u.input 2);
\draw[DEvod] (-.2,.55)|-(d.input 1);
\draw[DEvod] (d.input 2)--(-.5,0 |- d.input 2) node[left] {$A$};
\draw[DEvod] (u.output)--(2.8,1.1)|-(f.input 1);
\draw[DEvod] (d.output)--(2.8,-.1)|-(f.input 2);
\draw[DEvod] (f.output)--++(.4,0) node[right] {$F$};
\end{tikzpicture}

\end{frame}
